\documentclass[9pt,twocolumn]{extarticle}
\usepackage[utf8]{inputenc}
\usepackage[T1]{fontenc}
\usepackage[spanish]{babel}
\usepackage[a4paper,margin=1.25cm]{geometry}
\usepackage{microtype,graphicx,booktabs,tabularx}
\newcommand{\parencite}[1]{\textsuperscript{[#1]}}
\renewcommand{\cite}[1]{\textsuperscript{[#1]}}
\setlength{\parindent}{0pt}
\setlength{\parskip}{3pt}
\title{Notas de exposición: Netflix, RPC, gRPC y RMI}
\author{GA-SUM-02}
\date{21 de noviembre de 2026}
\begin{document}
\maketitle

\section{Investigación sobre la comunicación entre microservicios}

Para comprender por qué gRPC puede ser una alternativa útil en los sistemas distribuidos, se revisó el estudio de Niswar et al. (2024), en el que se compara el rendimiento de tres formas de comunicación: REST, GraphQL y gRPC.

Los investigadores realizaron sus pruebas en SISTER, un sistema de información educativa de Indonesia. Para el experimento desarrollaron tres microservicios utilizando Go, distribuidos en tres contenedores. El sistema incluía servicios de autenticación y consultas de información de docentes, además de utilizar Redis y MySQL para gestionar los datos.

El objetivo fue comparar el tiempo de respuesta y el uso del procesador de cada alternativa. Para ello, realizaron pruebas con diferentes cantidades de solicitudes y dos tipos de consultas: datos planos y datos anidados.

Uno de los resultados más importantes fue que gRPC obtuvo menores tiempos de respuesta en varias de las pruebas. Por ejemplo, al consultar datos planos con 100 solicitudes, registró un promedio de 233,84 milisegundos, mientras que REST alcanzó 1113,33 milisegundos y GraphQL 3852,07 milisegundos.

Sin embargo, los resultados también muestran que gRPC no fue superior en todos los aspectos. REST presentó un menor uso del procesador en las mediciones generales. Por esta razón, la elección de un protocolo debe considerar las necesidades y características del sistema.

Los autores relacionan el rendimiento de gRPC con el uso de HTTP/2, que permite enviar varias solicitudes por una misma conexión mediante la multiplexación. Esto facilita la comunicación cuando diferentes servicios necesitan intercambiar información al mismo tiempo.

Este estudio permite comprender algunas de las ventajas de gRPC en una arquitectura de microservicios. Aunque sus resultados corresponden al sistema SISTER y no a Netflix, sirven como referencia para analizar las características de este mecanismo de comunicación.

\subsection{Referencia consultada}

Niswar, M., Safruddin, R. A., Bustamin, A., y Aswad, I. (2024). Performance evaluation of microservices communication with REST, GraphQL, and gRPC. \textit{International Journal of Electronics and Telecommunications}, 70(2), 429--436. https://doi.org/10.24425/ijet.2024.149562

% Guion oral - Steven Pozo - 3 minutos
% Avance Día 1 - 20 de noviembre de 2026

\section*{Guion: arquitectura y modelo de comunicación}

\textbf{0:00--0:20. Alcance.}
Este diagrama no pretende mostrar toda la infraestructura de Netflix. Es una
vista lógica de un caso que Netflix sí documentó: su plataforma de comprensión
de medios. Así evitamos inventar nombres de servicios o conexiones internas.

\textbf{0:20--1:05. Componentes.}
A la izquierda separo dos consumidores. La aplicación o interfaz consumidora
entra por la ruta GraphQL documentada. El servicio backend consumidor es el
cliente gRPC: posee el stub y realiza la RPC unaria. Esta separación evita
afirmar, sin evidencia, que la aplicación final llama directamente a gRPC.
En el límite de la plataforma existen dos interfaces documentadas: GraphQL para
integraciones de interfaz y gRPC, usado principalmente de backend a backend.
La ruta que nos interesa llega al Search Gateway. Este módulo coordina el
procesamiento de la consulta y de la respuesta y se conecta con capacidades de
búsqueda e índices de datos de medios. Las cajas representan responsabilidades
lógicas, no máquinas específicas ni cantidades de réplicas.

\textbf{1:05--1:45. Cliente-servidor y RPC unaria.}
La aplicación es el cliente y el servicio gRPC es el servidor. Ambos comparten
un contrato definido con Protocol Buffers. El cliente llama un método en un
stub; gRPC serializa y transporta el mensaje por HTTP/2. En una RPC unaria el
cliente envía exactamente una petición y recibe una respuesta con datos o un
estado de error. Se parece a una función local, pero sigue siendo una llamada de
red y por eso puede tener latencia, vencimiento o fallo.

La secuencia se puede recordar con cuatro verbos: \textbf{invocar, serializar,
ejecutar y responder}. Primero se invoca el stub; luego se serializa la petición;
el servidor ejecuta la operación; finalmente devuelve la respuesta y el estado.
Si no llega a tiempo, el cliente no debe asumir que fue una llamada local exitosa.

\textbf{1:45--2:25. Por qué gRPC.}
La evidencia de Netflix indica dos necesidades: comunicación entre servicios
escritos en diferentes lenguajes y evitar bibliotecas cliente mantenidas a
mano. gRPC genera código desde el contrato y da una interfaz uniforme. En el
caso de la plataforma de medios, Netflix afirma explícitamente que gRPC se usa
principalmente entre backends; esa es la evidencia que sostiene la flecha del
diagrama.

\textbf{2:25--3:00. Límite y cierre.}
gRPC resuelve la invocación remota, pero no garantiza por sí solo que Netflix
sea disponible, consistente o tolerante a fallos. El diagrama no atribuye
replicación, balanceo, reintentos ni descubrimiento sin fuente. La idea clave es
separar el contrato de comunicación de la topología física y de los mecanismos
operativos.

\subsection*{Preguntas probables}

\begin{itemize}
  \item \textbf{¿Netflix usa gRPC en todo?} No. La fuente dice que predomina en
  backend a backend; la misma plataforma también expone GraphQL.
  \item \textbf{¿Una RPC es local?} No. El stub oculta parte de la mecánica,
  pero la llamada cruza la red y puede fallar.
  \item \textbf{¿Dónde entra RMI?} Como alternativa Java--Java para comparar;
  no hay evidencia presentada aquí de que Netflix use Java RMI.
  \item \textbf{¿Protocol Buffers y gRPC son lo mismo?} No. Protocol Buffers
  define y serializa mensajes; gRPC organiza la llamada remota y sus estados.
  \item \textbf{¿Por qué el diagrama no muestra regiones o réplicas?} Porque las
  fuentes utilizadas para este caso no demuestran esa topología física concreta.
\end{itemize}

\subsection*{Luis Quintero: CAP/PACELC y comparación (3 minutos)}

\textbf{0:00--0:55.} No voy a etiquetar a Netflix entero como AP o EL. Analizo
una operación concreta: lecturas locales y replicación entre regiones de
EVCache para datos no críticos. Netflix permite que la réplica remota se atrase
sin detener las operaciones locales. Ante una partición mantiene disponibilidad
y tolera temporalmente valores distintos entre regiones.

\textbf{0:55--1:35.} Sin partición, la misma decisión favorece latencia: se lee
localmente y se replica de forma asíncrona. Para ese caso limitado hablamos de
AP/EL. No depende de gRPC. Otros datos de Netflix usan cuórum y consistencia de
lectura después de escritura, así que una sola etiqueta no sirve para toda la
plataforma.

\textbf{1:35--2:30.} Como sistema cercano comparo la infraestructura de video
de YouTube/Google. Netflix publica una arquitectura de microservicios, cachés y
contratos gRPC. Google documentó software distribuido coordinado con
aceleradores VCU para transcodificar video y reportó entre 20 y 33 veces más
eficiencia. No afirmamos que YouTube use aquí gRPC porque la fuente no lo dice.

\textbf{2:30--3:00.} RPC es el modelo general; gRPC es un framework con
contratos protobuf y soporte multilenguaje; RMI es una alternativa centrada en
Java. La conclusión es que la tecnología de comunicación y la política de
consistencia son decisiones distintas.

% Incluir en las notas del equipo. Las citas respaldan el guion y no se leen.
\subsection*{Leonel Zambrano Triviño}
\textbf{Comunicación.} ¿Qué pasa si un servicio pide datos y reunirlos todos tarda demasiado? En Netflix hay una consulta llamada \texttt{GetItems}. El servidor puede devolver lo que ya encontró y una marca para pedir la siguiente página. Así, el cliente recibe datos a tiempo. Si ese cliente está atendiendo otra llamada con un plazo de gRPC, puede dejar de pedir páginas cuando se le acaba el tiempo \cite{leonel-kv,fallos-grpc}.

\smallskip
\textbf{Tolerancia a fallos.} Una llamada lenta y un servicio saturado se pueden parecer desde el celular. Para el cliente, el \emph{deadline} marca hasta cuándo esperar una respuesta. Para el servicio, una señal de saturación es el uso de CPU. Si sube demasiado, Netflix puede rechazar trabajo de baja prioridad para seguir atendiendo a quien acaba de pulsar reproducir \cite{fallos-grpc,fallos-netflix}.

Eso ocurrió en PlayAPI. Después de una caída, los equipos Android enviaron doce veces más solicitudes de precarga por segundo que antes. La precarga prepara contenido por adelantado. Su disponibilidad llegó a bajar al 20\,\%. En cambio, las solicitudes iniciadas por los usuarios mantuvieron una disponibilidad superior al 99,4\,\%. Ahí se ve qué protegió la decisión de dar prioridad \cite{fallos-netflix}.

La pelea entre Jake Paul y Mike Tyson muestra otro tipo de fallo. El 15 de noviembre de 2024, más de 60 millones de hogares siguieron la transmisión y miles reportaron pausas. ThousandEyes encontró errores durante el pico de demanda. Sus mediciones apuntan a congestión en los servidores que entregaban el video y no mostraron un problema general de los proveedores de internet. Una posible explicación es que miles de espectadores necesitaban el mismo segmento apenas se producía. La causa interna exacta queda abierta. En muchas pruebas, el video volvió al reconectar con una tasa de bits menor. El espectador recuperó la transmisión, aunque perdió calidad \cite{fallos-temple,fallos-thousandeyes}.

\smallskip
\textbf{Conclusión.} Para nuestro proyecto propongo que identifiquemos una acción que el usuario necesita completar y una tarea secundaria que pueda esperar. En una prueba de sobrecarga, mediríamos por separado cuántas peticiones de cada tipo terminan a tiempo. Si al posponer la tarea secundaria mejoran las respuestas de la acción principal, tendríamos evidencia para adoptar esa prioridad.




\section*{Declaración de uso de IA}
Se utilizó ChatGPT/Codex para organizar borradores, revisar claridad, proponer
una estructura de exposición y apoyar la compilación. El equipo revisó las
afirmaciones contra las fuentes citadas y conserva responsabilidad sobre el
contenido presentado.
\section*{Referencias principales}
Netflix Technology Blog (Media Understanding Platform, Service Mesh, EVCache,
capa clave--valor y load shedding); documentación oficial de gRPC; Niswar et
al. (2024), DOI 10.24425/ijet.2024.149562; Henríquez et al. (2025), DOI
10.15665/rp.v23i1.3683; y Ranganathan et al. (2021), DOI
10.1145/3445814.3446723. Los metadatos completos están en
\texttt{referencias/}.
\end{document}
